\section{Canonical approximation of the full pair function}\label{sec:full-pair}

The preceding source estimates were obtained on an off-axis disk with
a safe anchor. A full pair function must also retain real zeros and
zeros on both sides of the axis. Folding a transversely vanishing test
does not accomplish this. We instead use two reflected disks and a
smooth partition that counts every real label with total weight one.
This section contains the quantitative result \cite{OCGaussian},
accepted in the internal E07 assessment after the compact-center
qualification below; see Appendix~\ref{app:status}.

\subsection{A two-sided window with an explicit support buffer}

Define the $C^2$ cutoff
\begin{equation}\label{eq:full-cutoff}
 P(t)=\begin{cases}
 1,&t\le0,\\
 1-10t^3+15t^4-6t^5,&0<t<1,\\
 0,&t\ge1.
 \end{cases}
\end{equation}
Its first derivative is $-30t^2(1-t)^2$ on $(0,1)$, so
$|P'|\le15/8$; differentiation once more gives $|P''|\le6$.
For $X\ge10$ and $0<h\le1$ put
\begin{equation}\label{eq:full-window}
 \begin{gathered}
 w_{X,h}(x)=P(2|x-X|/h-1),\qquad \kappa_0(y)=P(|y|-1),\\
 W_X(x+\ii y)=w_{X,h}(x)\kappa_0(y),\qquad
 \theta_+(y)=P(1/2-16y),\quad\theta_-(y)=1-\theta_+(y).
 \end{gathered}
\end{equation}
The absolute values cause no singular derivatives because each cutoff
is constant near their nonsmooth points. The window equals one for
$|x-X|\le h/2$, $|y|\le1$, and vanishes outside
$|x-X|\le h$, $|y|\le2$. For true time-zero zeros the transverse
cutoff is one; for approximating zeros the collar $1<|y|<2$ keeps
its stated weights.

Let $c_\pm=X\pm4\ii$ and retain $R=9/2$, $\rho_*=13/3$, $r=25/6$
from \eqref{eq:radii}. The support of $\theta_+W_X$ has
$y\ge-1/32$ and hence
\begin{equation}\label{eq:two-disk-buffer}
 |z-c_+|^2\le1+(4+1/32)^2
        =\frac{17665}{1024}<\frac{625}{36}=r^2.
\end{equation}
The reflected statement holds for $\theta_-W_X$. Both supports have a
strict buffer inside their disks. At $y=0$, $\theta_+=\theta_-=1/2$.
Thus the partition retains real zeros exactly, without an axis indicator
or a boundary distribution.

Use $F_+=F=D H_0$ and $\widehat F_+=D\widehat H_{J,U}$ from
Section~\ref{sec:canonical}. Define
$F_-(z)=\overline{F_+(\bar z)}$ and
$\widehat F_-(z)=\overline{\widehat F_+(\bar z)}$.
Conjugation symmetry of both entire fields makes these products of the
same fields with reflected, locally zero-free multipliers. Reflection
transports the source error, anchor and log estimate to the lower disk:
\begin{equation}\label{eq:both-log-errors}
 \int_{D(c_\sigma,r)}
       |\log|F_\sigma|-\log|\widehat F_\sigma||\dd A\le E_X,
 \qquad \sigma\in\{+,-\}.
\end{equation}
Each field has at most $n_X$ zero labels in each $\rho_*$ disk, by
the same Jensen bound. Thus at most $M_X^{\rm pop}:=2n_X$ labels
of either field can have nonzero $W_X$ weight. This population bound
is distinct from the normalization constant $M_X=2X+26$.

\begin{lemma}[Two-sided weak transfer]\label{lem:two-sided}
For any compact $C^2$ test $\psi$ supported in the window rectangle,
\begin{equation}\label{eq:two-sided-transfer}
 \left|(\nu_{H_0}-\nu_{\widehat H_{J,U}})(\psi)\right|
 \le\frac{E_X}{2\pi}
       \sum_{\sigma=\pm}\|\Delta(\theta_\sigma\psi)\|_\infty.
\end{equation}
\end{lemma}
\begin{proof}
Write $\psi=\theta_+\psi+\theta_-\psi$. Each summand is supported
strictly within its inner disk. Apply \eqref{eq:weak-onepoint} there
to the difference of zero measures, remove each zero-free multiplier,
and use \eqref{eq:both-log-errors}. This argument treats all labels,
including coincident labels and labels on the axis, without simplicity
or zero separation assumptions.
\end{proof}

\subsection{Uniform Gaussian derivative costs}

All suprema in the following estimates are over the whole relevant
Euclidean space. Direct radial differentiation of \eqref{eq:gaussian}
gives
\begin{equation}\label{eq:gaussian-costs}
 \|G_\epsilon\|_\infty=A_\epsilon=\frac1{2\pi\epsilon^2},\quad
 \|\nabla G_\epsilon\|_\infty=B_\epsilon
       =\frac{A_\epsilon}{\sqrt{\ee}\epsilon},\quad
 \|\Delta G_\epsilon\|_\infty=C_\epsilon
       =\frac{2A_\epsilon}{\epsilon^2}.
\end{equation}
For the gradient, maximize $t\ee^{-t^2/2}$ at $t=1$; for the
Laplacian use $|t^2-2|\ee^{-t^2/2}\le2$ for $t\ge0$.

The derivative bounds for the window and partition are
$|w'|\le15/(4h)$, $|w''|\le24/h^2$,
$|\theta_\sigma'|\le30$, $|\theta_\sigma''|\le1536$.
The derivatives of $\theta_\sigma$ are supported in $|y|<1/32$,
where $\kappa_0=1$; the derivatives of $\kappa_0$ are supported in
$1<|y|<2$, where $\theta_\sigma$ is constant. This disjointness
removes their cross terms. With
$a_\sigma=\theta_\sigma W_X$ and
$b_\sigma=\theta_\sigma W_X^2$, we obtain
\begin{equation}\label{eq:window-costs}
 \begin{aligned}
 |a_\sigma|,|b_\sigma|&\le1,\\
 \|\nabla a_\sigma\|_\infty&\le L_h=15/(4h)+30,&
 \|\Delta a_\sigma\|_\infty&\le D_h=24/h^2+1536,\\
 \|\nabla b_\sigma\|_\infty&\le L_h^{(2)}=15/(2h)+30,&
 \|\Delta b_\sigma\|_\infty&\le D_h^{(2)}=609/(8h^2)+1536.
 \end{aligned}
\end{equation}
For example $|(w^2)''|\le2|w''|+2|w'|^2\le609/(8h^2)$.
Likewise $|(\kappa_0^2)''|\le609/32<1536$, so the vertical
bound is dominated by the partition transition. Bounding the Euclidean
gradient by the sum of its coordinate bounds proves the stated $L$'s.

For brevity write $A=A_\epsilon$, $B=B_\epsilon$, $C=C_\epsilon$,
and define
\begin{equation}\label{eq:UV-costs}
 \begin{split}
 U_{h,\epsilon}&=D_hA+2L_hB+C,\\
 V_{h,\epsilon}&=D_h^{(2)}A^2+4L_h^{(2)}AB+2AC+2B^2.
 \end{split}
\end{equation}
The product rule gives, uniformly in $u,v\in\C$,
\begin{align}
 \|\Delta_z[a_\sigma(z)G_\epsilon(u-z)]\|_\infty
       &\le U_{h,\epsilon},\label{eq:single-test-cost}\\
 \|\Delta_z[b_\sigma(z)G_\epsilon(u-z)G_\epsilon(v-z)]\|_\infty
       &\le V_{h,\epsilon}.\label{eq:self-test-cost}
\end{align}
For the second line, the Gaussian product has sup norm at most $A^2$,
gradient norm at most $2AB$, and Laplacian norm at most $2AC+2B^2$.

\subsection{The full pair approximation theorem}

Use the same $W_X$ and $\epsilon$ for the true field and its exact
finite-source approximant, and let $R_{2,\epsilon}^{W_X}(V)$ denote
\eqref{eq:full-correlation} formed from the zeros of $V$.

\begin{theorem}[Uniform full-pair source approximation]\label{thm:full-pair}
Let $X\ge10$, $0<h\le1$, $\epsilon>0$, and suppose the complete
source error \eqref{eq:complete-delta} obeys
$\delta\le\min(1,c_0/2)$. With the explicit constants of
\eqref{eq:canonical-constants}, \eqref{eq:EX}, and
\eqref{eq:UV-costs},
\begin{equation}\label{eq:full-pair-error}
 \sup_{u,v\in\C}
 \left|R_{2,\epsilon}^{W_X}(H_0)(u,v)
          -R_{2,\epsilon}^{W_X}(\widehat H_{J,U})(u,v)\right|
 \le\frac{E_X}{\pi}
       \left(4n_XA_\epsilon U_{h,\epsilon}+V_{h,\epsilon}\right).
\end{equation}
The estimate includes every real or nonreal zero label in the window,
with its multiplicity and its actual cutoff weight.
\end{theorem}
\begin{proof}
Applying Lemma~\ref{lem:two-sided} to
$W_X(z)G_\epsilon(u-z)$ gives
\[
 \sup_u|r(u)-\widehat r(u)|\le E_XU_{h,\epsilon}/\pi.
\]
The population bound gives $0\le r,\widehat r\le2n_XA_\epsilon$.
Thus
\[
 |r(u)r(v)-\widehat r(u)\widehat r(v)|
 \le\frac{E_X}{\pi}\,4n_XA_\epsilon U_{h,\epsilon},
\]
by adding and subtracting $\widehat r(u)r(v)$. Applying the same
lemma to $W_X(z)^2G_\epsilon(u-z)G_\epsilon(v-z)$ bounds
$|d(u,v)-\widehat d(u,v)|$ by $E_XV_{h,\epsilon}/\pi$.
Subtract the self terms using Proposition~\ref{prop:weighted-pair}.
All estimates are uniform in the observation points; their Gaussian
tails do not alter the compact support of the tests in $z$.
\end{proof}

For fixed $h,X$ and $0<\epsilon\le1$, the right side is
$O_{h,X}(E_X\epsilon^{-6})$. This follows from
$A=O(\epsilon^{-2})$, $B=O(\epsilon^{-3})$,
$C=O(\epsilon^{-4})$, so $AU$ and $V$ are both
$O_h(\epsilon^{-6})$. The estimate is conservative; no optimal
width exponent is claimed. Nor does a uniform density bound imply
an $L^1$ bound on the whole noncompact observation space without
additional tail estimates.

Set $K_{X,h,\epsilon}=(4n_XA_\epsilon U_{h,\epsilon}
+V_{h,\epsilon})/\pi$. For any density tolerance $e>0$, choose
\begin{equation}\label{eq:full-precision}
 \delta_* =\min\left\{1,\frac{c_0}{2},
 M_X\left(\frac{e}{2K_{X,h,\epsilon}C_X}\right)^{1/\alpha_X}\right\}.
\end{equation}
Choose the arithmetic and radial tails in \eqref{eq:delta-H} each
at most $\delta_*/(2\mathcal B_X)$. Their explicit decay gives finite
$J,U$, and \eqref{eq:full-pair-error} is then at most $e/2$.
The remaining $e/2$ is available to a future certified evaluator.
No evaluation of the canonical pair function, or efficient runtime,
is asserted by this existence argument.

\subsection{Center averaging and the remaining source covariance}

The uniformity in observation points makes the theorem compatible
with \eqref{eq:unfolding}. If each center has an error bound
$\mathcal E_X$ from \eqref{eq:full-pair-error}, then
\begin{equation}\label{eq:averaged-transfer}
 \sup_{\zeta,\eta}
   |\overline R_2(\zeta,\eta)-\widehat{\overline R}_2(\zeta,\eta)|
       \le\int_I J_X^2\mathcal E_X\dd\pi(X).
\end{equation}
This follows by the triangle inequality. For common finite source cutoffs
at a prescribed uniform density tolerance, use the full-pair windows
\eqref{eq:full-window} with $0<h_{\min}\le h_X\le1$, together with
the positive lower bound on $\epsilon_X$ already declared. Then
$K_{X,h_X,\epsilon_X}$ and $J_X^2K_{X,h_X,\epsilon_X}$ are bounded
on the compact center interval. Finite source cutoffs can be chosen uniformly: the positive anchor is common, $M_X,\mathcal B_X,n_X,C_X$
are bounded, and $\alpha_X$ has a positive lower bound. This is a
compact-center assertion, not an all-height uniform estimate.

There is also an exact source expression for the first intensity.
Define the compact kernels
\[
 K_{X,\sigma,\zeta}(z)=
  \Delta_z\{\theta_\sigma(\Im z)W_X(z)
                   G_{\epsilon_X}(T_X\zeta-z)\}.
\]
The partition and zero-measure identity give
\begin{equation}\label{eq:source-intensity}
 \widetilde r_X(\zeta)
   =\frac{J_X}{2\pi}\sum_{\sigma=\pm}
       \int_{D(c_\sigma,r)}\log|F_{X,\sigma}(z)|
                                  K_{X,\sigma,\zeta}(z)\dd A(z).
\end{equation}
The analogous single-log expression for $\widetilde d_X$ uses
$J_X^2\Delta_z(\theta_\sigma W_X^2G_{\epsilon_X}(T_X\zeta-z)
G_{\epsilon_X}(T_X\eta-z))/(2\pi)$.
Inserting these formulas into \eqref{eq:center-covariance} identifies
a centered product of localized log integrals. Its kernels, center
law and diagonal subtraction are explicit. The individual log bounds
control approximation error; they do not establish cancellation or a
limiting shape for this center covariance. That source-specific joint
estimate is the principal next analytical question.
